\thispagestyle{fancy}
\begin{center}
\textbf{2026 HiMCM --- Problem A}\\[9pt]
{\Large\bfseries Choosing a City for High-School Life\\with Incomplete Evidence}\\[12pt]
\textbf{Summary Sheet}
\end{center}
A city recommendation becomes more useful when its evidence and priorities can
be examined together. We compare sixteen cities using ten indicators of safety,
environment, household costs, education and everyday resources. Zurich is the
conservative recommendation under the baseline and equal weights. All sixteen
cities remain possible champions within the retained evidence bounds.

The evaluation combines fixed-anchor TOPSIS with a maximum-regret decision.
Benefit and cost utilities share fixed reference scales, allowing controlled
comparisons across data, preference and development scenarios. Missing values
remain intervals. We select the city with the smallest worst-case shortfall from
a better-scoring competitor over the declared rectangular uncertainty set.

Official municipal records, school statistics and operator inventories supplement
seventeen previously missing entries while preserving the team's original values.
The average score-interval width falls from 18.81077 to \CurrentMeanWidth{}, a \TotalWidthReduction\%
reduction. The resulting 160-entry matrix contains 112 point values, 29 partial
intervals and 19 fully missing intervals. Actual reference years and source
populations remain explicit.

A controlled formulation revision retains the 48 uncertain entries rather
than collapsing each to a midpoint. Although Zurich remains the recommendation, \RevisionRankChanges{}
fixed-weight decision ranks change. A completion strictly inside the intervals
lets Dublin exceed Zurich, establishing why a midpoint first place cannot
certify a certain champion.

Preferences and development conditions answer different decision questions.
Doubling the relative emphasis on education changes the conservative choice to
Oslo. \SummaryWeightMC{} These are conditional simulation frequencies.
Three uniform ten-per-cent development scenarios retain Zurich; separate
city-specific shocks reveal conditional recommendation changes without claiming
an empirical forecast.

Separate country, district and age-group cases examine what citywide totals
leave unresolved. They distinguish facility supply from student access and
observed broad age groups from unsupported age-specific preferences. The result
is a practical starting choice with explicit alternatives and clear priorities
for additional evidence.

\textbf{Keywords:} high-school opportunities; TOPSIS; interval evidence;
maximum regret; preference sensitivity.
\clearpage
\tableofcontents
\clearpage
\section{Introduction and decision scope}
\subsection{What makes a city suitable for high-school life?}
A useful comparison connects a student's school day with the opportunities and
constraints around it. Safety and air quality affect the daily setting; teaching
resources and nearby higher education describe educational supply; housing,
living costs, transport, travel, broadband and swimming facilities describe
household pressures and everyday resources. Our operational definition of
``best'' is the conservative choice under these ten criteria, declared weights
and retained evidence ranges. It provides a starting recommendation for the
compared cities, rather than a claim about every student's experience.

The study uses the sixteen cities in the frozen team matrix: Amsterdam, Berlin,
Dublin, Helsinki, London, Madrid, Milan, Oslo, Paris, Prague, Singapore,
Stockholm, Tokyo, Vienna, Warsaw and Zurich. This is an inherited comparison
set across sixteen countries, not a representative sample of all world cities.
It exceeds the problem's minimum of ten cities and spans different resource
and cost settings. The findings concern this set; a broader claim would require
an explicit sampling design and additional comparable observations.

\subsection{Criterion selection and contribution}
The ten indicators preserve the agreed team framework and balance education
with the conditions that make resources usable in daily life. Their directions
and anchors appear in Table~\ref{tab:indicators}. Counts and city averages serve
as scoped supply or setting measures. They do not directly measure friendships,
student wellbeing, school quality, language fit, discrimination, cultural
participation or individual access. Those outcomes are relevant to a family's
choice; comparable student-level observations are needed to integrate them
without inventing measurements or double-counting existing proxies.

Our central contribution is a city decision that retains incomplete evidence.
We improve source precision, compare a point-completion benchmark with an
interval revision, and examine preferences and future conditions without
changing the primary indicator set. TOPSIS provides one main quality score;
AHP structures stated priorities, Monte Carlo measures specified perturbations,
and VIKOR supplies a focused method comparison. The extensions test the
recommendation rather than vote among an expanding set of ranking models.

AI assistance supported research retrieval, calculations, code and prose
preparation, including Exa source retrieval and extraction \cite{openai,exa}.
The appended use report records the known assistance
and the outstanding documentation required before submission.

\section{Assumptions and their role}
\begin{enumerate}
\item \textbf{Comparability follows the retained scope.} Original values and
population denominators remain fixed. Each supplement keeps its provider,
year, city boundary and educational population. Mixed-year densities are
conditional comparisons; they do not create a same-year census.
\item \textbf{Utilities are monotone on fixed scales.} More resources and less
cost or burden improve their respective utilities. Declared anchors set the
scale; they are neither estimated preferences nor empirical optimum values.
\item \textbf{The score permits compensation.} Weighted distances combine
criteria in a common utility space. The model therefore allows stronger results
on some criteria to offset weaker results on others. A family requiring a hard
safety or eligibility threshold should apply it separately before choosing.
\item \textbf{The evidence set is rectangular.} Uncertain utilities vary
independently within retained intervals. This conservative closure omits unknown
cross-criterion and cross-city restrictions. Its possible champions and regret
bounds are conditional on this set; they are not population probabilities.
\item \textbf{Extensions hold the comparison basis fixed.} Scenario changes
retain population denominators, anchors and the indicated weight rule. They
are conditional experiments with no fitted future date or representative
age-specific preference distribution.
\end{enumerate}

\section{Notation}
\begin{tabular}{p{.18\linewidth}p{.75\linewidth}}
\toprule Symbol & Meaning\\\midrule
$M,N$ & Sixteen cities and ten indicators\\
$w_j,a_j$ & Normalized criterion weight and fixed utility anchor\\
$u_{ij}$ & Utility of indicator $j$ for city $i$, on $[0,1]$\\
$\ell_{ij},h_{ij}$ & Retained lower and upper utility bounds\\
$L_i,U_i$ & Lower and upper TOPSIS score, on $[0,100]$\\
$R_i,\mathcal C$ & Maximum regret and set of possible champions\\
$\kappa$ & Concentration of a declared Dirichlet weight distribution\\
$r$ & Raw-value multiplier in a development scenario\\\bottomrule
\end{tabular}
